\begin{theorem}
Let $s\to\infty$ through positive integers and let $a\geq0$ be an integer
sequence with $a=o(s)$.  In the standard eventual sense, for every
$\varepsilon>0$ and all sufficiently large indices,
\[
 r(s,s+a)\geq(1-\varepsilon)\frac{s}{e}
  2^{(s+a-1)/2-a^2/(2s)}.
\]
Equivalently, this is
\[
 r(s,s+a)\geq(1+o(1))\frac{s}{e}
 2^{(s+a-1)/2-a^2/(2s)}.
\]
\end{theorem}
\begin{proof}
Take $s$ and $a$ along the given sequences and put $k=s+a$.  Eventually
$4\leq s\leq k$, so \noderef{F2ForwardIndependentBound} gives a loopless
$T_s$-free digraph on
\[
 N=2^{2s-3}-2^{s-1}-2^{s-2}+1=(1+o(1))2^{2s-3}
\]
vertices and
\[
 \operatorname{fwi}_k(D)\leq
 2^{3s^2/2+as-5s/2+o(s)}.
\]
The $o(s)$ estimate follows directly from its rank sum: writing $t=s-j$, use
$\binom{s+a}{j+a}\leq(e(s+a)/(a+1))^{j+a}$, maximize the resulting quadratic
in $j$, and use $a=o(s)$ to obtain
$(e(s+a)/(a+1))^a=2^{o(s)}$.

Apply \noderef{RandomPermutationReduction}.  The resulting $K_s$-free graph
$G$ has
\[
 i_k(G)\leq(e/k)^k
 2^{3s^2/2+as-5s/2+o(s)}.
\]
If this count is zero, then $r(s,k)\geq N$, which is stronger than the
claimed estimate.  Otherwise choose $p=i_k(G)^{-1/k}$.  This lies in
$[0,1]$ and satisfies the hypothesis of \noderef{SamplingDeletion}, so
\[
 r(s,k)>N i_k(G)^{-1/k}-1.
\]
Substitute the preceding count and $N=(1+o(1))2^{2s-3}$.  Since $k=s+a$,
the exponent simplifies as
\[
 2s-3-\frac{3s^2/2+as-5s/2+o(s)}{s+a}
 =\frac{s+a-1}{2}-\frac{a^2}{2s}+o(1),
\]
where $a=o(s)$ makes the discarded terms $o(1)$ after division by $s+a$.
The factor $k/e$ is $(1+o(1))s/e$, and the $-1$ is negligible.  Thus the
eventual $(1-\varepsilon)$ inequality follows for every $\varepsilon>0$.
\end{proof}
